% Proposed isolated A.89. No frozen29 source is changed.
\statement{A.89}{Proved}{The actual hexagon defect and the full-GT dyadic coefficient comparison}
Let \(\alpha=(-1,f)\in\widehat{GT}_0\) be the actual negative involution,
with canonical \(u\in F'\), and retain the full marked data of
A.87--A.88. Put
\[
 \begin{gathered}
 z=(xy)^{-1},\qquad d=d_\alpha,\qquad m=(d-1)/2\in2\widehat{\mathbb Z},\\
 u_0=u(x,y),\quad u_1=u(y,z),\quad u_2=u(z,x),\\
 \mathcal H_d=u_0x^m u_2z^m u_1y^m .
 \end{gathered}
 \tag{A.89.1}
\]
This is the entire actual hexagon word, whether or not it is trivial.
Let \(K=\overline{\langle\!\langle x\rangle\!\rangle}\),
\(\Lambda=\mathbb F_2[[\widehat{\mathbb Z}]]\), and identify
\(K/\overline{K^2[K,K]}=\Lambda[x]\), with retained head \(T=y\)
and \([x]=1\), as in A.88. For the involution on coefficients write
\(\overline{P}(T)=P(T^{-1})\), and set
\[
 U=[u_0],\qquad W=[u_1],\qquad R_d=[\mathcal H_d],\qquad
 \Theta_u=1+(T-1)(U+\overline U).
 \tag{A.89.2}
\]
Use A.88's full continuous geometric sums
\((T-1)D_b=T^b-1\). Then
\[
 \boxed{\ \overline{R_d}+T^dR_d=\Theta_uD_d+1\ },
 \qquad
 \boxed{\ a_{\rm reg}=d_\alpha
       \ \Longleftrightarrow\ \overline{R_d}=T^dR_d\ }.
 \tag{A.89.3}
\]
Under the additional hypothesis \(\alpha\in\widehat{GT}\), one has
the entire dyadic equality
\[
                       a_{{\rm reg},2}=d_{\alpha,2}.
 \tag{A.89.4}
\]
Neither the full coefficient equality nor the nonlinear
\(\Xi_d=\Omega_\alpha=1\) is asserted.

\begin{proof}
\emph{1. Two actual basis substitutions.}\par
The free basis \((y,z)\) has \(x=(yz)^{-1}\). Its swap \(\beta\) satisfies
\[
 \beta(y)=z,\qquad \beta(z)=y,\qquad
 \beta(x)=(zy)^{-1}=y^{-1}xy.
\]
It preserves K, inverts its retained head, and sends its module
generator to \(T^{-1}[x]\). Its full module action is therefore
\(P(T)\mapsto T^{-1}P(T^{-1})\). Canonical duality in this same
basis gives \(\beta(u_1)=u(z,y)=u_1^{-1}\). In characteristic two,
\[
                         \overline W=TW.
 \tag{A.89.5}
\]
The weight is the actual conjugator \(y^{-1}xy\); it cannot be dropped.
Next the genuine basis substitution \(x\mapsto x,\ y\mapsto z\)
preserves K, fixes \([x]\), and inverts the retained head. It acts
by coefficient bar. Canonical duality consequently gives
\[
                  [u_2]=[u(z,x)]=\overline U.
 \tag{A.89.6}
\]
All three words are substitutions of the same full canonical u,
not independent module parameters.

\emph{2. Ordered evaluation of the entire word.}\par
Use kernel-first coordinates
\[
 (P,T^a)(Q,T^b)=(P+T^aQ,T^{a+b}).
\]
The actual \(z=y^{-1}x^{-1}\) has module coordinate \(T^{-1}\)
and head \(T^{-1}\). Hence for every profinite m,
\[
                 z^m=(T^{-m}D_m,T^{-m}).
 \tag{A.89.7}
\]
The formula follows on integer powers from geometric sums and
then by continuity. The module coordinate of \(x^m\) is zero
because m is even. Reading (A.89.1) in its stated order gives
\[
 R_d=U+\overline U+T^{-m}D_m+T^{-m}W.
 \tag{A.89.8}
\]
Only the later \(u_1\) term carries the accumulated head \(T^{-m}\).
No group factors have been commuted.

The full geometric-sum identities are
\[
 \overline{D_m}=T^{1-m}D_m,\qquad
 1+T^m=(T-1)D_m,\qquad
 D_{2m+1}=1+T(1+T^m)D_m.
 \tag{A.89.9}
\]
They follow first on integer powers and then by the already
established continuity of \(D_b\); no division in a finite group
algebra is used. Put \(F_u=U+\overline U\), which is bar-invariant.
Equations (A.89.5)--(A.89.9), with \(d=2m+1\), give
\[
 \begin{aligned}
 \overline{R_d}+T^dR_d
  &=(1+T^d)F_u+T(1+T^m)D_m
        +T^m\overline W+T^{m+1}W\\
  &=(T-1)D_dF_u+D_d+1\\
  &=\Theta_uD_d+1.
 \end{aligned}
 \tag{A.89.10}
\]
The W terms cancel by their actual weighted relation.
A.88 proves \(\Theta_uD_{a_{\rm reg}}=1\) and its exact
full-head coefficient test, so (A.89.3) follows.

\emph{3. The additional full-GT consequence.}\par
Now assume the full pentagon hypothesis
\(\alpha\in\widehat{GT}\). A.42 applies to this very canonical
u, and supplies its prescribed pro-two positive pair for either
local square root of \(1+24c_2(u)_2\). The complete A.45 comparison
identifies the roots with \(\pm d_{\alpha,2}\) and already
gives \(d_{\alpha,2}\equiv1\pmod4\). Thus the prescribed pair
\((d_{\alpha,2},u_2)\) satisfies its actual pro-two positive hexagon.
The image of the particular word \(\mathcal H_d\) in the
maximal pro-two quotient of F is consequently one.

The compatible map of the actual cyclic-cover module to
\(\Lambda_2=\mathbb F_2[[\mathbb Z_2]]\) therefore sends
\(R_d\) to zero. Projecting (A.89.3) and A.88 gives
\[
       \Theta_{u,2}D_{d_{\alpha,2}}
          =1=\Theta_{u,2}D_{a_{{\rm reg},2}}.
 \tag{A.89.11}
\]
The first equality makes \(\Theta_{u,2}\) a unit. Cancel it
and multiply by \(T-1\) to obtain
\(T^{d_{\alpha,2}}=T^{a_{{\rm reg},2}}\).
Finite cyclic two-power heads separate \(\mathbb Z_2\);
their group elements are distinct basis elements in the finite
group algebra. This proves (A.89.4). Equivalently, the two
canonical frame coordinates of A.87 satisfy
\(t_2=\rho_{D,2}\), and the dyadic coordinate of
\((\Omega_\alpha)_{\rm ab}\) vanishes.

The pentagon input in this last step is indispensable: it is
not supplied for an arbitrary \(\widehat{GT}_0\)-involution.
No uniqueness of a projected norm frame or relative fixed
subgroup has been assumed. Even the full coefficient identity
would determine only this Alexander quotient, not the entire
group word \(\Xi_d\). The prospective assignment
\(s\mapsto v,\ \delta\mapsto\delta^d\) is still not a modular
homomorphism before its cubic relation is proved; no relative
lifting theorem is applied to an unestablished input.
\end{proof}
